\section{Working backward and working forward}\label{ch04:sec:directions}
Every proof has a \emph{goal}, the statement it must establish, and a supply of facts it may use: the hypotheses, the definitions, and earlier results. When you do not know what to write next, look in two directions. \emph{Working backward} starts from the goal and asks: what would be enough to establish this? \emph{Working forward} starts from the facts and asks: what follows from what I already have? A proof is found when the two meet: something derived from the facts is exactly what working backward said would be enough.

Unpacking the definitions in the goal often tells you exactly what to produce. To show that an integer is even, you must find an integer $w$ such that the number equals $2w$ (Definition~\ref{def:parity}). To show that a function $f$ is injective, you take inputs $x$ and $y$ with $f(x)=f(y)$ and deduce that $x=y$ (Section~\ref{ch03:sec:injsurj}). To show that an object with some property exists, you usually name one and check that it has the property; such an object is called a \emph{witness} (Section~\ref{ch01:sec:definition}). Occasionally existence is proved without a witness, by showing that if no such object existed, something impossible would follow (a proof by contradiction, Section~\ref{ch04:sec:contrapositive}); Chapter~\ref{ch:probability} contains proofs of this kind.

\subsection*{A proof state on paper}
Consider the claim
\begin{center}
\emph{If $n$ is an even integer, then $n^2$ is even.}
\end{center}
Section~\ref{ch02:sec:implication} proved it in one line; here we slow down to watch how such a proof is found. At each moment of the search we keep two lists: the facts we already have, and what we still have to show. Together they are called the \emph{proof state}. Proof assistants such as Lean display exactly this pair after every step (Chapter~\ref{ch:lean1}).

At the start, the only fact is the hypothesis, and the goal is ``$n^2$ is even.'' Working backward, the definition of even turns the goal into \emph{find an integer $w$ with $n^2=2w$}. Working forward, the same definition turns the hypothesis into $n=2k$ for some integer $k$. The two letters play different roles: $k$ is handed to us by the hypothesis, while $w$ is an integer we still have to find. Squaring gives $n^2=4k^2=2(2k^2)$, exactly the shape the goal asks for, so we choose $w=2k^2$, an integer because $k$ is. Only now does the goal move to the ``known'' list. The table records the search.
\begin{center}
\begin{tabularx}{\linewidth}{@{}l>{\raggedright\arraybackslash}X>{\raggedright\arraybackslash}X@{}}
\toprule
Step & Known & Still to show\\
\midrule
Start & $n$ is an even integer & $n^2$ is even\\
Unpack the goal & $n$ is an even integer & an integer $w$ with $n^2=2w$\\
Unpack the hypothesis & $n=2k$, where $k$ is an integer & an integer $w$ with $n^2=2w$\\
Square and factor & $n=2k$ and $n^2=2(2k^2)$ & an integer $w$ with $n^2=2w$\\
Choose $w=2k^2$ & $n^2=2w$, and $w$ is an integer & nothing: the proof is finished\\
\bottomrule
\end{tabularx}
\end{center}
The search moved in both directions, but the written proof runs forward only:
\begin{quote}
Let $n$ be an even integer. Then $n=2k$ for some integer $k$, so $n^2=4k^2=2(2k^2)$. Since $2k^2$ is an integer, $n^2$ is even.
\end{quote}
A proof like this one, which starts from the hypotheses and reaches the conclusion through a chain of deductions, is called a \emph{direct proof}.

Keeping the goal on the ``still to show'' list until it has been earned guards against a common mistake: writing ``since $n^2$ is even, $n^2=2w$ for some integer $w$'' in the middle of the proof. That uses the conclusion as though it were already known, so the argument goes in a circle. Working backward is still the right way to \emph{search}, but in the finished proof the goal may appear only at the end, as the conclusion reached.

\pauseq{Let $f:\ZZ\to\ZZ$ be given by $f(n)=n+3$. To prove that $f$ is surjective, what does the goal ask you to produce? Find it by working backward, and then check it by working forward.}{By the definition of surjective (Section~\ref{ch03:sec:injsurj}), for each integer $b$ you must find an integer $a$ with $a+3=b$; working backward, this suggests $a=b-3$. Working forward, $b-3$ is an integer and $f(b-3)=(b-3)+3=b$. So $b-3$ is a witness for each $b$, and $f$ is surjective.}
